\documentclass[10pt,a4paper]{amsart}
\usepackage{amsmath,amssymb,amsthm,graphicx}
\usepackage{verbatim}
\newenvironment{reference}{\comment}{\endcomment}
%\newenvironment{reference}{}{}
\newenvironment{slogan}{\comment}{\endcomment}
\newenvironment{history}{\comment}{\endcomment}

% For commutative diagrams we use Xy-pic
\usepackage[all]{xy}

% We use 2cell for 2-commutative diagrams.
\xyoption{2cell}
% Theorem environments.
%
\theoremstyle{plain}
\newtheorem{theorem}[subsection]{Theorem}
\newtheorem{proposition}[subsection]{Proposition}
\newtheorem{lemma}[subsection]{Lemma}

\theoremstyle{definition}
\newtheorem{definition}[subsection]{Definition}
\newtheorem{example}[subsection]{Example}
\newtheorem{exercise}[subsection]{Exercise}
\newtheorem{situation}[subsection]{Situation}

\theoremstyle{remark}
\newtheorem{remark}[subsection]{Remark}
\newtheorem{remarks}[subsection]{Remarks}

\numberwithin{equation}{subsection}

% Macros
%
\def\lim{\mathop{\mathrm{lim}}\nolimits}
\def\colim{\mathop{\mathrm{colim}}\nolimits}
\def\Spec{\mathop{\mathrm{Spec}}}
\def\Hom{\mathop{\mathrm{Hom}}\nolimits}
\def\Ext{\mathop{\mathrm{Ext}}\nolimits}
\def\SheafHom{\mathop{\mathcal{H}\!\mathit{om}}\nolimits}
\def\SheafExt{\mathop{\mathcal{E}\!\mathit{xt}}\nolimits}
\def\Sch{\mathit{Sch}}
\def\Mor{\mathop{\mathrm{Mor}}\nolimits}
\def\Ob{\mathop{\mathrm{Ob}}\nolimits}
\def\Sh{\mathop{\mathit{Sh}}\nolimits}
\def\NL{\mathop{N\!L}\nolimits}
\def\CH{\mathop{\mathrm{CH}}\nolimits}
\def\proetale{{pro\text{-}\acute{e}tale}}
\def\etale{{\acute{e}tale}}
\def\QCoh{\mathit{QCoh}}
\def\Ker{\mathop{\mathrm{Ker}}}
\def\Im{\mathop{\mathrm{Im}}}
\def\Coker{\mathop{\mathrm{Coker}}}
\def\Coim{\mathop{\mathrm{Coim}}}

% Boxtimes
%
\DeclareMathSymbol{\boxtimes}{\mathbin}{AMSa}{"02}

%
% Macros for moduli stacks/spaces
%
\def\QCohstack{\mathcal{QC}\!\mathit{oh}}
\def\Cohstack{\mathcal{C}\!\mathit{oh}}
\def\Spacesstack{\mathcal{S}\!\mathit{paces}}
\def\Quotfunctor{\mathrm{Quot}}
\def\Hilbfunctor{\mathrm{Hilb}}
\def\Curvesstack{\mathcal{C}\!\mathit{urves}}
\def\Polarizedstack{\mathcal{P}\!\mathit{olarized}}
\def\Complexesstack{\mathcal{C}\!\mathit{omplexes}}
% \Pic is the operator that assigns to X its picard group, usage \Pic(X)
% \Picardstack_{X/B} denotes the Picard stack of X over B
% \Picardfunctor_{X/B} denotes the Picard functor of X over B
\def\Pic{\mathop{\mathrm{Pic}}\nolimits}
\def\Picardstack{\mathcal{P}\!\mathit{ic}}
\def\Picardfunctor{\mathrm{Pic}}
\def\Deformationcategory{\mathcal{D}\!\mathit{ef}}
\usepackage{cleveref}
\crefname{originalreference}{original reference}{original references}
\Crefname{originalreference}{Original reference}{Original references}
\makeatletter
\newlabel{situation-shriek}{{0F42}{source}{}{sourcepointer-1}{}}
\newlabel{situation-shriek@cref}{{[originalreference][0][]0F42}{[1][0][]source}{}{}{}}
\makeatother
\crefname{originaldefinition}{original definition}{original definitions}
\Crefname{originaldefinition}{Original definition}{Original definitions}
\makeatletter
\newlabel{more-morphisms-definition-CM}{{045R}{source}{}{sourcepointer-2}{}}
\newlabel{more-morphisms-definition-CM@cref}{{[originaldefinition][0][]045R}{[1][0][]source}{}{}{}}
\makeatother
\crefname{originallemma}{original lemma}{original lemmas}
\Crefname{originallemma}{Original lemma}{Original lemmas}
\makeatletter
\newlabel{morphisms-lemma-flat-finite-presentation-CM-fibres-relative-dimension}{{02NM}{source}{}{sourcepointer-3}{}}
\newlabel{morphisms-lemma-flat-finite-presentation-CM-fibres-relative-dimension@cref}{{[originallemma][0][]02NM}{[1][0][]source}{}{}{}}
\makeatother
\crefname{originallemma}{original lemma}{original lemmas}
\Crefname{originallemma}{Original lemma}{Original lemmas}
\makeatletter
\newlabel{lemma-shriek-coherent}{{0AU1}{source}{}{sourcepointer-4}{}}
\newlabel{lemma-shriek-coherent@cref}{{[originallemma][0][]0AU1}{[1][0][]source}{}{}{}}
\makeatother
\crefname{originallemma}{original lemma}{original lemmas}
\Crefname{originallemma}{Original lemma}{Original lemmas}
\makeatletter
\newlabel{lemma-flat-shriek}{{0BV6}{source}{}{sourcepointer-5}{}}
\newlabel{lemma-flat-shriek@cref}{{[originallemma][0][]0BV6}{[1][0][]source}{}{}{}}
\makeatother
\crefname{originallemma}{original lemma}{original lemmas}
\Crefname{originallemma}{Original lemma}{Original lemmas}
\makeatletter
\newlabel{lemma-flat-quasi-finite-shriek}{{0BV7}{source}{}{sourcepointer-6}{}}
\newlabel{lemma-flat-quasi-finite-shriek@cref}{{[originallemma][0][]0BV7}{[1][0][]source}{}{}{}}
\makeatother
\crefname{originallemma}{original lemma}{original lemmas}
\Crefname{originallemma}{Original lemma}{Original lemmas}
\makeatletter
\newlabel{more-morphisms-lemma-flat-finite-presentation-characterize-CM}{{0BUU}{source}{}{sourcepointer-7}{}}
\newlabel{more-morphisms-lemma-flat-finite-presentation-characterize-CM@cref}{{[originallemma][0][]0BUU}{[1][0][]source}{}{}{}}
\makeatother
\crefname{originallemma}{original lemma}{original lemmas}
\Crefname{originallemma}{Original lemma}{Original lemmas}
\makeatletter
\newlabel{lemma-shriek-affine-line}{{0AA1}{source}{}{sourcepointer-8}{}}
\newlabel{lemma-shriek-affine-line@cref}{{[originallemma][0][]0AA1}{[1][0][]source}{}{}{}}
\makeatother
\title{For a Cohen-Macaulay morphism the shriek of the structure sheaf is a flat module in one degree}
\author{The Stacks Project Authors}
\date{stacks-project:duality.tex:local-snapshot-20261012; GNU FDL 1.2 or later}
\begin{document}
\maketitle
\noindent\textit{Editorial note.} All delivered mathematical text is the author's own, copied from the pinned Stacks Project chapter file duality.tex; no statement or proof step was rewritten, reordered or removed.\par\medskip
\noindent\textit{Editorial note.} Preamble provenance: the pinned chapter file opens with an input of the shared preamble and carries no preamble text of its own, so the project's single Stacks-wide preamble staged at sources/stacks-duality/preamble.tex is excerpted verbatim under the recipe fixed by LANE\_BRIEF 22.1, as three spans, lines 8-13 (verbatim package and the reference/slogan/history environments), lines 14-18 (Xy-pic setup) and lines 154-217 (the native theorem declarations and mathematical macros). Nothing was added, removed or redefined; the wrapper class and the amsmath/amssymb/amsthm/graphicx base packages are supplied by the assembler.\par\medskip
\noindent\textit{Editorial note.} The counted statement and proof contain no cross reference, no citation and no included file, so the delivered body is self contained.\par\medskip
\noindent\textit{Editorial note.} Canonical location of the counted proposition: Stacks tag 0BV8 (sources/stacks-duality/pinned-original-source.tex line 5978).\par\medskip
\noindent\textit{Allowed standard prerequisites (editorial).}\begin{itemize}
\item The relative dimension of a Cohen-Macaulay morphism is locally constant. Morphisms, Lemma labelled morphisms-lemma-flat-finite-presentation-CM-fibres-relative-dimension, cited at proof lines 6015-6017. Fixes d.
\item Coherence of the shriek of the structure sheaf and flatness of its single cohomology sheaf. Same chapter, Lemmas labelled lemma-shriek-coherent and lemma-flat-shriek, cited at proof lines 6017-6025. Give coherence and flatness once vanishing is known.
\item The affine line computation and the quasi-finite case. Same chapter, Lemmas labelled lemma-shriek-affine-line and lemma-flat-quasi-finite-shriek, cited at proof lines 6033-6039. Settle the two model cases of the reduction.
\item Characterisation of Cohen-Macaulay morphisms through factorization. More on Morphisms, Lemma labelled more-morphisms-lemma-flat-finite-presentation-characterize-CM, cited at proof lines 6029-6033. Produces the factorization with g quasi-finite flat.
\end{itemize}\medskip
\ifdefined\hypertarget\hypertarget{sourcepointer-1}{}\fi
\noindent\textit{Source reference.} Original reference 0F42: sources/stacks-duality/pinned-original-source.tex, source lines 3953--3953. Same-chapter pointer into the author environment situation at source line 3953; invoked at proof line None of the counted proof. It is retained only as a navigation pointer; every load-bearing step is inside the delivered span.\par\medskip
\ifdefined\hypertarget\hypertarget{sourcepointer-2}{}\fi
\noindent\textit{Source reference.} Original definition 045R: sources/stacks-chow/tags.txt, source lines 5446--5446. Cross-chapter standard prerequisite located in the project tag list staged in the local snapshot; invoked at proof line None of the counted proof. It is retained only as a navigation pointer; every load-bearing step is inside the delivered span.\par\medskip
\ifdefined\hypertarget\hypertarget{sourcepointer-3}{}\fi
\noindent\textit{Source reference.} Original lemma 02NM: sources/stacks-chow/tags.txt, source lines 3485--3485. Cross-chapter standard prerequisite located in the project tag list staged in the local snapshot; invoked at proof line 5993 of the counted proof. It is retained only as a navigation pointer; every load-bearing step is inside the delivered span.\par\medskip
\ifdefined\hypertarget\hypertarget{sourcepointer-4}{}\fi
\noindent\textit{Source reference.} Original lemma 0AU1: sources/stacks-duality/pinned-original-source.tex, source lines 4686--4686. Same-chapter pointer into the author environment lemma at source line 4686; invoked at proof line 5995 of the counted proof. It is retained only as a navigation pointer; every load-bearing step is inside the delivered span.\par\medskip
\ifdefined\hypertarget\hypertarget{sourcepointer-5}{}\fi
\noindent\textit{Source reference.} Original lemma 0BV6: sources/stacks-duality/pinned-original-source.tex, source lines 5915--5915. Same-chapter pointer into the author environment lemma at source line 5915; invoked at proof line 5996 of the counted proof. It is retained only as a navigation pointer; every load-bearing step is inside the delivered span.\par\medskip
\ifdefined\hypertarget\hypertarget{sourcepointer-6}{}\fi
\noindent\textit{Source reference.} Original lemma 0BV7: sources/stacks-duality/pinned-original-source.tex, source lines 5962--5962. Same-chapter pointer into the author environment lemma at source line 5962; invoked at proof line 6003 of the counted proof. It is retained only as a navigation pointer; every load-bearing step is inside the delivered span.\par\medskip
\ifdefined\hypertarget\hypertarget{sourcepointer-7}{}\fi
\noindent\textit{Source reference.} Original lemma 0BUU: sources/stacks-chow/tags.txt, source lines 14887--14887. Cross-chapter standard prerequisite located in the project tag list staged in the local snapshot; invoked at proof line 6009 of the counted proof. It is retained only as a navigation pointer; every load-bearing step is inside the delivered span.\par\medskip
\ifdefined\hypertarget\hypertarget{sourcepointer-8}{}\fi
\noindent\textit{Source reference.} Original lemma 0AA1: sources/stacks-duality/pinned-original-source.tex, source lines 4609--4609. Same-chapter pointer into the author environment lemma at source line 4609; invoked at proof line 6010 of the counted proof. It is retained only as a navigation pointer; every load-bearing step is inside the delivered span.\par\medskip
\begin{lemma}
\label{lemma-CM-shriek}
In Situation \ref{situation-shriek} let $f : X \to Y$ be a morphism of
$\textit{FTS}_S$. If $f$ is Cohen-Macaulay (More on Morphisms, Definition
\ref{more-morphisms-definition-CM}), then
$$
f^!\mathcal{O}_Y = \omega_{X/Y}[d]
$$
for some coherent $\mathcal{O}_X$-module $\omega_{X/Y}$ flat over $Y$
where $d$ is the locally constant
function on $X$ which gives the relative dimension of $X$ over $Y$.
\end{lemma}

\begin{proof}
The relative dimension $d$ is well defined and locally constant by
Morphisms, Lemma
\ref{morphisms-lemma-flat-finite-presentation-CM-fibres-relative-dimension}.
The cohomology sheaves of $f^!\mathcal{O}_Y$ are coherent by
Lemma \ref{lemma-shriek-coherent}.
We will get flatness of $\omega_{X/Y}$ from Lemma \ref{lemma-flat-shriek}
if we can show the other cohomology sheaves of $f^!\mathcal{O}_Y$
are zero.

\medskip\noindent
The question is local on $X$, hence we may assume $X$ and $Y$ are affine
and the morphism has relative dimension $d$. If $d = 0$, then the
result follows directly from Lemma \ref{lemma-flat-quasi-finite-shriek}.
If $d > 0$, then we may assume there is a factorization
$$
X \xrightarrow{g} \mathbf{A}^d_Y \xrightarrow{p} Y
$$
with $g$ quasi-finite and flat, see More on Morphisms, Lemma
\ref{more-morphisms-lemma-flat-finite-presentation-characterize-CM}.
Then $f^! = g^! \circ p^!$. By Lemma \ref{lemma-shriek-affine-line}
we see that $p^!\mathcal{O}_Y \cong \mathcal{O}_{\mathbf{A}^d_Y}[-d]$.
We conclude by the case $d = 0$.
\end{proof}

\end{document}
